\section{Contexto y problema}

\subsection{Cómo se realiza hoy el proceso}
En la Facultad de Ciencias de la Computación de la UTEQ, la selección de docentes no
titulares se realiza sin un sistema propio. Según el Decano y la Coordinadora de la
carrera, los postulantes envían su hoja de vida por correo a Talento Humano, que la
reenvía al decanato, o directamente al decanato. Desde allí se descargan y se revisan
una por una.

El Decano y la Coordinadora califican cada hoja de vida con una matriz de valoración
en Excel, que asigna puntaje a los títulos, la experiencia docente, las publicaciones
afines a la cátedra, los cursos de actualización y los reconocimientos académicos. Los
mejores candidatos son llamados a una entrevista, cuyo puntaje se suma al de la matriz.
Como requisito mínimo, el postulante debe tener título de cuarto nivel y experiencia
docente.

Con los puntajes se elabora un informe y un acta, que se envían mediante oficio al
Consejo Directivo de la facultad. Este lo aprueba y lo traslada al Consejo Académico,
y luego al Consejo Universitario, que es el órgano colegiado superior. Como este
consejo a veces se reúne solo una vez al mes, la rectora puede aprobar la contratación
de forma administrativa y el consejo la ratifica después.

\subsection{Problemas identificados}
\begin{itemize}
  \item Según el Decano, la etapa que más tiempo consume es revisar las hojas de vida
        y registrar los puntajes en la matriz.
  \item El informe se redacta y los puntajes se suman a mano. La Coordinadora señaló
        que el proceso \enquote{tiene que automatizarse}.
  \item Las hojas de vida llegan por distintos canales de correo, sin un
        repositorio único ni seguimiento del estado de cada postulante.
  \item En las entrevistas no se obtuvieron cifras de tiempo por etapa ni del número
        de postulantes por convocatoria; se registran como pendientes de confirmar.
\end{itemize}

\subsection{Interesados}
\begin{table}[H]
\centering
\caption{Interesados y expectativas}
\begin{tabularx}{\linewidth}{@{}p{4cm}X@{}}
\toprule
\textbf{Interesado} & \textbf{Qué espera del sistema} \\
\midrule
Decano & Almacenar las hojas de vida y generar la matriz de valoración, que considera la parte más complicada. \\
Coordinadora de carrera & Automatizar el cálculo de puntajes y la elaboración del informe. \\
Postulante & Enviar sus documentos y consultar el estado de su postulación. \\
Comisión evaluadora & Administrar las entrevistas desde la aplicación. \\
Autoridades (consejos y rectorado) & Recibir el informe con los puntajes para aprobar la contratación. \\
\bottomrule
\end{tabularx}
\end{table}